\section{Results}
\label{sec:results}

\subsection{Final comparison}

The untouched base scored \PaperScores{0/12}{8/8}{8/8}. The adapter selected
from \Code{checkpoint-84} scored \PaperScores{11/12}{8/8}{8/8}. Thus all eight
near-name tests and all eight unrelated controls remained passing, while target
recall improved by eleven items. Every tuned output was non-empty, so the
canonical decision passed and the evidence records the result as
\Code{acceptance-approved} with interpretation
\Code{candidate-knowledge-acquisition}.\claimsource{manifest}
\claimsource{evaluation}

The sole tuned recall failure, record \Code{fact\_006}, returned ``I do not
know.'' It was a miss rather than an incorrect transfer to another entity. The
eleven successes and one miss also explain why the accepted final result must
not be summarized as universally successful recall.

\begin{table}[htbp]
\centering
\caption{Fixed 28-row final suite. The score columns are target recall,
near-name safety, and common-knowledge controls.}
\label{tab:final}
\begin{tabular}{lccc}
\toprule
System & Recall & Safety & Controls \\
\midrule
Untouched pinned base & 0/12 & 8/8 & 8/8 \\
Selected BF16 LoRA & 11/12 & 8/8 & 8/8 \\
\bottomrule
\end{tabular}
\end{table}

\subsection{Checkpoint trajectory}

Table~\ref{tab:trajectory} reports every saved evaluation. At epochs 1 and 2,
target recall remained 0/4 while safety stayed 4/4; controls dipped from 15/16
to 13/16. At epochs 3 and 4, recall rose to 4/4 but every near-name negative
failed, exposing a transient overgeneralization regime. The first 24/24
checkpoint appeared at epoch 5. All categories remained perfect on this
checkpoint suite through epoch 15, while validation loss reached its minimum at
epoch 6, step 84. That tie-break selected \Code{checkpoint-84} even though the
run completed the full 210/210 steps and 15 epochs.\claimsource{evaluation}

\begin{longtable}{rrrrrrl}
\caption{All fifteen checkpoint-suite evaluations. Loss is printed to nine
decimal places from the machine record.}\label{tab:trajectory}\\
\toprule
Epoch & Step & Recall & Safety & Controls & Eval. loss & Selection \\
\midrule
\endfirsthead
\toprule
Epoch & Step & Recall & Safety & Controls & Eval. loss & Selection \\
\midrule
\endhead
\TrajectoryRow{1}{14}{0/4}{4/4}{15/16}{0.849535763}{--}
\TrajectoryRow{2}{28}{0/4}{4/4}{13/16}{0.275242716}{--}
\TrajectoryRow{3}{42}{4/4}{0/4}{16/16}{0.092094362}{--}
\TrajectoryRow{4}{56}{4/4}{0/4}{16/16}{0.143437117}{--}
\TrajectoryRow{5}{70}{4/4}{4/4}{16/16}{0.049332991}{--}
\TrajectoryRow{6}{84}{4/4}{4/4}{16/16}{0.040961619}{selected}
\TrajectoryRow{7}{98}{4/4}{4/4}{16/16}{0.059368324}{--}
\TrajectoryRow{8}{112}{4/4}{4/4}{16/16}{0.064691097}{--}
\TrajectoryRow{9}{126}{4/4}{4/4}{16/16}{0.065394856}{--}
\TrajectoryRow{10}{140}{4/4}{4/4}{16/16}{0.067192160}{--}
\TrajectoryRow{11}{154}{4/4}{4/4}{16/16}{0.066379592}{--}
\TrajectoryRow{12}{168}{4/4}{4/4}{16/16}{0.068035968}{--}
\TrajectoryRow{13}{182}{4/4}{4/4}{16/16}{0.068136089}{--}
\TrajectoryRow{14}{196}{4/4}{4/4}{16/16}{0.066940054}{--}
\TrajectoryRow{15}{210}{4/4}{4/4}{16/16}{0.067972727}{--}
\bottomrule
\end{longtable}

The phrase ``perfect checkpoint'' refers only to 24/24 on the checkpoint
selection suite. The disjoint final suite was larger and produced 11/12 recall.
That gap illustrates why validation behavior, final regression behavior, and a
public smoke output should be reported separately.

\subsection{Time, memory, and cost}

The user-facing invocation lasted 2,106 seconds (35 minutes, 6 seconds). The
trainer reported 1,584.9314 seconds and 0.132 optimizer steps per second.
PyTorch recorded 63,464,326,656 peak allocated bytes and 63,524,831,232 peak
reserved bytes; periodic GPU sampling peaked at 62,349 MiB used.
\claimsource{metadata}

The command-window estimate was \$0.83 at the observed active rate. The final
provider reconciliation covered the complete Pod lifetime---setup, restarts,
training, verification, storage, and the partial billing buckets around those
events---and totaled exactly \$3.2853100409265606. We report that as \$3.29,
not as the narrower command estimate.\claimsource{billing}
